\specialchapter{Preface}

This work gathers together, without substantial modification, the majority of the historical Notes which have appeared to date in my Eléments de Mathématique. Only the flow has been made independent of the Eléments to which these Notes were attached; they are therefore, in principle, accessible to every reader who possesses a sound classical mathematical background, of undergraduate standard.

Of course, the separate studies which make up this volume could not in any way pretend to sketch, even in a summary manner, a complete and connected history of the development of Mathematics up to our day. Entire parts of classical mathematics such as differential Geometry, algebraic Geometry, the Calculus of variations, are only mentioned in passing; others, such as the theory of analytic functions, that of differential equations or partial differential equations, are hardly touched on; all the more do these gaps become more numerous and more important as the modern era is reached. It goes without saying that this is not a case of intentional omission; it is simply due to the fact that the corresponding chapters of the Eléments have not yet been published.

Finally the reader will find in these Notes practically no bibliographic or anecdotal information about the mathematicians in question; what has been attempted above all, for each theory, is to bring out as clearly as possible what were the guiding ideas, and how these ideas developed and reacted the ones on the others.

The numbers in square brackets refer to the Bibliography at the end of the volume.

\tableofcontents
\mainmatter

\chapter{FOUNDATIONS OF MATHEMATICS; LOGIC; SET THEORY.}

The study of what is usually called the ``foundations of mathematics'', which has been carried out ceaselessly since the beginning of the 19th century, was impossible to bring to fruition except with the help of a parallel effort to systematise Logic, at least those parts that govern the links between mathematical statements. So it is not possible to separate the history of set theory and the formalisation of mathematics from that of ``Mathematical Logic''. But traditional logic, like that of the modern philosophers, covers in principle, an area of applications far greater than Mathematics. Therefore the reader must not expect to find in what follows a history of Logic, even in a very summarised form; we have limited ourselves as far as possible to retracing the evolution of Logic only in so far as it impinged on that of Mathematics. It is because of this that we will say nothing about the non-classical logics (many-valued logics, modal logics); all the more so will we be unable to tackle the history of those controversies which, from the Sophists to the Vienna School, have never stopped dividing philosophers both as to the possibility and the manner of applying Logic to objects in the real world or to concepts of human thought.

That there was a well-developed prehellinic mathematics is not today in any doubt. Not only are the notions (already very abstract) of whole number and of the measurement of quantities commonly used in the most ancient documents which have reached us from Egypt or Chaldea, but Babylonian algebra, because of the elegance and sureness of its methods, should not be thought of as a simple collection of problems solved by empirical fumbling. And, if nothing is found in the texts which resembles a ``proof'' in the formal meaning of the word, it is reasonable to believe that the discovery of such methods of solution, whose generality appears through particular numerical applications, was not possible without a minimum of logical links (perhaps not entirely conscious, but rather like those on which a modern algebraist depends when he undertakes a calculation, before ``setting down formally'' all the details) ({[}232{]}, pp.~203 ff.).

The essential originality of the Greeks consisted precisely of a conscious effort to order mathematical proofs in a sequence such that passing from one link to the next leaves no room for doubt and constrains universal assent.

That Greek mathematicians made use, in the course of their research, just like modern mathematicians, of ``heuristic'', rather than convincing, arguments, is what was proved for example (if there were any need), by the ``treatise of method'' of Archimedes {[}153 c{]}; note also in this, allusions to results ``found, but not proven'' by earlier mathematicians. \(^{1}\) But from the first detailed texts that are known to us (and which date from the middle of the Vth century), the ideal ``canon'' of a mathematical text is properly settled. It will find its highest expression in the great classics, Euclid, Archimedes and Apollonius; the notion of proof, in these authors, differs in no way from ours.

We have no texts allowing us to follow the first steps of this ``deductive method'', which seems to us already near perfection at the exact moment when we become aware of its existence. One can only think that it fits fairly naturally into the perpetual search for ``explanations'' of the world which characterises Greek thought and which is so discernible already amongst the Ionian philosophers of the VIIth century; moreover tradition is unanimous in ascribing the development and refinement of the method to the Pythagorean School, in a period which fits between the end of the VIth century and the middle of the Vth century.

It is on this ``deductive'' mathematics, fully conscious of its goals and methods, that the philosophical and mathematical thought of subsequent times will be concentrated. We will see on the one hand the establishment little by little of ``formal'' Logic modelled on mathematics, to conclude with the creation of formalised languages; on the other hand, mainly starting at the beginning of the XIXth century, the basic concepts of mathematics will be queried more and more and a great effort will be made to clarify their nature, especially after the birth of the Theory of sets.

\section*{THE FORMALISATION OF LOGIC.}
\phantomsection
\addcontentsline{toc}{section}{THE FORMALISATION OF LOGIC.}

The general impression which seems to emerge from the (very patchy) texts that we have on Greek philosophical thought of the Vth century, is that it is dominated by an increasingly conscious effort to extend to the whole field of human thought the procedures for conducting discussion put in hand with so much success by contemporary rhetoric and mathematics - in other words, to create Logic in its most general form. The tone of philosophical writings undergoes at this time a sudden change: while in the VIIth and VIth centuries the philosophers affirm or vaticinate (or at most outline vague arguments, based on equally vague analogies), starting with Parmenides and especially Zenon, they argue and try to draw out general principles which can serve as a basis for their dialectic: it is in Parmenides that one finds the first statement of the principle of the excluded middle, and the proofs ``by contradiction'' of Zenon of Elea remain famous. But Zenon writes in the middle of the Vth century; and, whatever the uncertainties of our documentation, \(^{2}\) it is very likely that at this time, mathematicians, in their own sphere, were currently using these principles.

As we said above, it is not for us to retrace the innumerable difficulties which abound at every step in the gestation of this Logic, and the controversies which result, from the Eleatics to Plato and Aristotle, via the Sophists; let us pick up here only the role played in this evolution by the assiduous cultivation of oratorical art and the analysis of language which is one of its corollaries, developments that it is agreed should be attributed mainly to the Sophists of the Vth century. On the other hand, even if the influence of mathematics is not always explicitly recognised, it is none the less manifest, in particular in the writings of Plato and Aristotle. It can be said that Plato was almost obsessed by mathematics; without being himself an inventor in this area, he kept up, after a particular period in his life, with the discoveries by contemporary mathematicians (many of whom were his friends or his pupils), and never after ceased from a most direct kind of interest, even going so far as to suggest new directions of research; so it is that constantly, in his writings, mathematics serves as illustration or model (and sometimes even, as with the Pythagoreans, feeding his leanings towards mysticism). As for his pupil Aristotle, he could not have avoided receiving the minimum of mathematical foundations required of the pupils of the Academy, and a volume has been produced consisting of extracts of his work that are about mathematics or refer to it {[}153 d{]}; but he seems never to have made a great effort to keep in touch with the mathematical movements of his time, and he only quotes in this area results that had been popularised a long time before. This displacement will besides only become more marked with the majority of later philosophers, of whom many, due to the lack of technical preparation, believe themselves in all good faith to be speaking knowledgeably about mathematics, whereas they will only be referring to a stage long since overtaken in its evolution.

The end product of this period, as far as Logic is concerned, is the monumental work of Aristotle {[}6{]}, whose great merit is to have succeeded in systematising and codifying for the first time procedures of reasoning which had remained vague or unformulated in his predecessors. \(^{3}\) We must above all retain here, as our aim, the general thesis of this work, to know that it is possible to reduce all correct reasoning to the systematic application of a small number of immutable rules, which are independent of the particular nature of the objects in question (an independence which is clearly demonstrated by the representation of the concepts or propositions by means of letters --- probably borrowed by Aristotle from the mathematicians). But Aristotle concentrates his attention almost exclusively on a particular type of relations and logical chains, making up what he calls the ``syllogism'': it is mainly a case of relations that we would translate nowadays in the form \(A \subset B\) or \(A \cap B \neq \emptyset\) in the language of the theory of sets, \(^{4}\) and with the way of chaining these relations or their negations, by means of the schema

\[
(A \subset B \text {   and   } B \subset C) \Rightarrow (A \subset C).
\]

Aristotle was still sufficiently informed of the mathematics of his era in order to have seen that schemas of this type were not sufficient to take account of all the logical operations used by mathematicians, nor, all the more, of the other applications of Logic ({[}6{]}, An. Pr., I, 46). \(^{5}\) At least the deep study of the different forms of ``syllogisms'' which he undertakes (and which is almost entirely consecrated to the elucidation of the perpetual difficulties raised by the ambiguity or the obscurity of the objects on which the reasoning bears) gives him among others the chance to formulate the rules for obtaining the negation of a proposition ({[}6{]}, An. Pr., I, 46). It is also to Aristotle that is due the credit for having distinguished with great precision the role of ``universal'' propositions from that of ``particular'' propositions, the first sketch of quantifiers. \(^{6}\) But it is too well known how the influence of his writings (often interpreted in a narrow and unintelligent way), which remain still very influential until well into the XIXth century, was to encourage philosophers in their neglect of the study of mathematics, and to block the progress of formal Logic. \(^{7}\)

However this latter continued to make progress in Antiquity, in the surroundings of the Megaric and Stoic schools, rivals of the Peripatetics. Our information about these tenets are unfortunately all at second hand, often passed on by adversaries or mediocre commentators. The essential progress achieved by these logicians consists, it would appear, of the formation of a ``propositional calculus'' in the meaning which it has today: instead of being restricted, like Aristotle, to propositions of the particular form \(A \subset B\) , they state rules about completely indeterminate propositions. Further, they had analysed the logical links between these rules in such a deep way that they knew how to deduce them all from five of them, set down as ``unprovable'', by means of procedures very similar to modern methods {[}23{]}. Unfortunately their influence was fairly ephemeral, and their results were to fall into oblivion until the day when they were rediscovered by the logicians of the XIXth century. Aristotle remains as the uncontested master in Logic until the XVIIth century; it is known in particular that the scholastic philosophers are entirely under his sway, and if their contribution to formal logic is far from negligible {[}25{]}, it does not contain any progress at the highest level compared with the achievements of the philosophers of Antiquity.

It is appropriate, however, to note here that it does not seem that the works of Aristotle or his successors had a notable influence on mathematics. The Greek mathematicians conducted their research along the path opened by the Pythagoreans and their successors of the IVth century (Theodorus, Theeta, Eudoxus) without apparently bothering with formal logic in the presentation of their results: a finding that should hardly astonish when one compares the flexibility and precision acquired already by that time by mathematical reasoning, to the very rudimentary state of Aristotelian logic. And when logic will overtake this stage, it is again the new acquisitions from mathematics that will guide it in its evolution.

With the development of algebra, it was indeed impossible to avoid being struck by the analogy between the rules of formal logic and the rules of algebra, the ones like the others having in common the property of being applicable to objects (propositions or numbers) which are not precisely determined. And when in the XVIIth century algebraic notation took its definitive form in the hands of Viète and Descartes, almost immediately one can see appearing diverse attempts at a symbolic notation intended for the representation of logical operations; but before Leibniz, these tentative efforts, for example like that of Herigone (1644) at writing down the proofs of elementary Geometry, or that of Pell (1659) at writing down those of Arithmetic, remain very superficial and do not lead to any progress in the analysis of mathematical reasoning.

With Leibniz, we are in the presence of a philosopher who is also a mathematician of the first rank, and who will know how to draw out of his mathematical experience the germ of the ideas that will bring formal logic out of the scholastic dead-end. \(^{8}\) A universal spirit if ever there were one, inexhaustible source of original and fruitful ideas, Leibniz would be interested even more in Logic as it found itself at the very heart of his grand projects for formalising language and thought, at which he never ceased working throughout his life. Trained in his childhood in scholastic logic, he was seduced by the idea (going back to Raymond Lulle) of a method which would reduce all human concepts to primitive concepts, making up an ``Alphabet of human thought'', and would rebuild them in a quasi mechanical way to obtain all true propositions ({[}198 b{]}, v. VII, p.~185; cf.~{[}69a{]}, chap.~II). Still very young, he had also conceived another much more original idea, that of the usefulness of symbolic notations as an ``Ariadne's string'' string of thought: \(^{9}\) ``The true method'', he said, ``must provide us with a filum Ariadnes, that is to say a kind of sensitive and coarse means that guides the mind, in the same way as lines drawn in geometry and the type of operations that are prescribed to apprentices in Arithmetic. Without that our mind would not know how to go along a long path without straying.'' ({[}198 b{]}, v. VII, p.~22; cf.~{[}69 a{]}, p.~90). Knowing little of the mathematics of his time until about his 25th year, it is at first in the form of a ``universal language'' that he presents his projects ({[}69 a{]}, chap.~III); but as soon as he comes in contact with Algebra, he adopts it as model for his ``universal Characteristic''. By that he means a type of symbolic language, capable of expressing without ambiguity all human thought, of reinforcing our powers of reasoning, of avoiding errors by means of an effort of entirely mechanical concentration, finally constructed so that ``the chimera that the person himself who presents them does not hear, would be impossible to write down in these characters'' ({[}198 a{]}, v. I, p.~187). In innumerable passages in his writings where Leibniz alludes to this grandiose project and to the progress to which its attainment would lead (cf.~{[}69 a{]}, chap.~IV and VI), it is seen with what clarity he had conceived the notion of a formalised language, a pure combination of symbols of which only the chaining is important, \(^{10}\) in such a way that a machine would be able to produce all the theorems, \(^{11}\) and that all controversies would be resolved by a simple calculation ({[}198 b{]}, v. VII, p.~198-203). If these hopes might appear excessive, it is none the less true that it is to this constant tendency of Leibniz's thought that must be tied a good part of his mathematical work, beginning with his work on the symbolism of the infinitesimal Calculus (see p.~190 ff.); he was himself perfectly aware of this, and linked explicitly also to his ``Characteristic'' his ideas on indicial notation and determinants ({[}198 a{]}, v. II, p.~204; cf.~{[}69 a{]}, pp.~481-487) and his sketch of ``geometrical Calculus'' (see pp.~50 and 61 ff.; cf.~{[}69 a{]}, chap IX). But in his mind the essential part had to be symbolic Logic, or, as he puts it a ``Calculus ratiocinator'', and though he does not succeed in creating this calculus, at least we see him trying to at least three times. On his first attempt, he has the idea of associating with each ``primitive'' term a prime number, each term made up of several primitive terms being represented by the product of the corresponding prime numbers; \(^{12}\) he seeks to translate into this system the usual rules of syllogism, but runs up against major complications caused by negation (that he tries, fairly naturally, to represent by a change of sign) and abandons rapidly this path ({[}198 c{]}, pp.~42-96; cf.~{[}69 a{]}, pp.326-344). In later attempts, he seeks to give Aristotelian logic a more algebraic form; sometimes he keeps the notation \(AB\) for the conjunction of two concepts; sometimes he uses the notation \(A + B\) ; \(^{13}\) he notes (in multiplicative notation) the law of idempotence \(AA = A\) , remarks that one can substitute the proposition ``all A is B'' by the equality \(A = AB\) and that one can recover starting from there most of the rules of Aristotle by a purely algebraic calculation ({[}198 c{]}, pp.~229-237 and 356-399; cf.~{[}69 a{]}, pp.345-364); he also has the idea of the empty set (``non Ens''), and recognises for instance the equivalence of the propositions ``all A is B'' and ``A.(not B) is not'' (loc. cit.). Further, he remarks that his logical calculus is applicable not only to the logic of concepts, but also to that of propositions ({[}198 c{]}, p.~377). He appears to be therefore very near to ``Boolean calculus''. Unfortunately, it appears that he was not able to free himself completely from the scholastic influence; not only does he set himself as almost the only aim of his calculus the transcription, into his notation, of the rules of syllogism, \(^{14}\) but he goes as far as sacrificing his most felicitous ideas to the desire to recover completely the rules of Aristotle, even those that were incompatible with the idea of an empty set. \(^{15}\)

The work of Leibniz remained for the most part unpublished until the beginning of the XXth century, and had only little direct influence. During the whole of the XVIIIth and the beginning of the XIXth centuries, different authors (de Segner, J. Lambert, Ploucquet, Holland, De Castillon, Gregonne) sketched attempts similar to those of Leibniz, without ever progressing substantially beyond the point at which he had stopped; their work only had a feeble effect, which means that most of them ignored all results due to their predecessors. \(^{16}\) It is in any case under the same conditions that G. Boole, who must be considered to be the real creator of modern symbolic logic {[}29{]}, writes. His key idea consists in putting himself systematically in the position of considering ``extension'', so of calculating directly with sets, and writing xy for the intersection of two sets, and \(x + y\) for their union when x and y have no element in common. He introduces as well a ``universe'' denoted 1 (the set of all elements) and the empty set denoted 0, and he writes 1 - x for the complement of x. As Leibniz had done, he interprets the relation of inclusion by the relation xy = x (from which he extracts without difficulty the justification of the rules of classical syllogism) and his notation for union and complements gives his system a flexibility that was missing in his forerunners. \(^{17}\) Furthermore, by associating with each proposition the set of ``cases'' in which it holds, he interprets the relation of implication as an inclusion, and his calculus of sets gives him in this way the rules of the ``propositional calculus''.

In the second half of the XIXth century, Boole's system formed the basis for the work of an active school of logicians who improved it and completed it at several points. It is thus that Jevons (1864) enlarged the meaning of the operation of union \(x + y\) by extending it to the case in which \(x\) and \(y\) are arbitrary; A. de Morgan in 1858 and C. S. Peirce in 1867 proved the duality laws

\[
(\mathbf {C} A) \cap (\mathbf {C} B) = \mathbf {C} (A \cup B), (\mathbf {C} A) \cup (\mathbf {C} B) = \mathbf {C} (A \cap B); ^{18}
\]

De Morgan tackles also, in 1860, the study of relations, defining inversion and the composition of binary relations (that is to say the operations corresponding to the operations \(\overline{G}^{-1}\) and \(G_{1} \circ G_{2}\) on graphs). \(^{19}\) All this work is systematically expounded and developed in the massive and prolific work of Schröder {[}277{]}. But it is fairly interesting to note that the logicians of whom we have been speaking do not seem very interested in the application of their results to mathematics and that, on the contrary, Boole and Schröder especially seem to have as their principal aim to develop ``Boolean'' algebra by imitating the methods and problems of classical algebra (often in a fairly artificial way). The reasons for this attitude must doubtless be seen in the fact that the Boolean calculus still lacked a facility for transcribing most of mathematical reasoning, \(^{20}\) and only supplied in this way a very partial answer to the great dream of Leibniz. The construction of a formalism better adapted to mathematics --- of which the introduction of variables and quantifiers, due independently to Frege {[}117 a, b, c{]} and C. S. Peirce {[}248 b{]}, make up the major stage --- was the work of logicians and mathematicians who, unlike the above, had above all as their aim applications to the foundations of mathematics.

Frege's project {[}117 b and c{]} was to create a foundation for arithmetic based on a logic formalised by a ``writing of concepts'' (Begriffschrift) and we will come back later (p.~29) to the way in which he defines the natural numbers. His work is characterised by extreme precision and attention to detail in the analysis of concepts; it is because of this tendency that he introduces many a distinction which turns out to be of great importance in modern logic: for instance it is he who first distinguishes between the statement of a proposition and the assertion that this proposition is true, between the relation of belonging and that of inclusion, between an object x and the set \(\{x\}\) reduced to this single object, etc.. His formalised logic, which contains not only ``variables'' with the meaning used in mathematics, but also ``propositional variables'' representing indeterminate relations, susceptible to quantification, would later (through the work of Russell and Whitehead) supply the fundamental tool for metamathematics. Unfortunately, the symbols which he adopts are hardly inspiring, of a terrifying typographic complexity and far removed from the practice of mathematicians; which had the result of turning away these latter and reducing considerably the influence of Frege on his contemporaries.

Peano's goal was at the same time much vaster and much more down to earth; it consisted in publishing a ``Formulary of mathematics'', written entirely in formal language and containing, not only mathematical logic, but all the results of the most important branches of mathematics. The speed with which he managed to complete this ambitious project, helped by a host of enthusiastic collaborators (Vailati, Pieri, Padoa, Vacca, Vivanti, Fano, Burali-Forti) is witness to the excellence of the symbolism which he had adopted: following closely the current practice of mathematicians, and introducing numerous well-chosen abbreviating symbols, his language remains as well fairly easily legible, thanks notably to an ingenious system for replacing brackets by full stop separators {[}246 f{]}. Much notation due to Peano is today adopted by the majority of mathematicians: we quote \(\in, \supset\) (but, contrary to the present use, with the meaning of ``is contained in'' or ``implies'' \(^{21}\) ), \(\cup, \cap, A - B\) (set of differences \(a - b\) where \(a \in A\) and \(b \in B\) ). On the other hand, it is in the ``Formulary'' that is found for the first time a thorough analysis of the general notion of function, of those of direct image \(^{22}\) and reciprocal image, and the remark that a sequence is only a function defined on N. But quantification, with Peano, is subject to hampering restrictions (one can only, in principle, quantify, in his system, relations of the form \(A \to B\) , \(A \Leftrightarrow B\) or \(A = B\) ). Further the almost fanatical zeal of some of his disciples laid them wide open to ridicule; criticism, often unjustified, by H. Poincaré in particular, was a heavy blow to the Peano school and became an obstacle to the diffusion of his doctrines in the mathematical world.

With Frege and Peano the essential elements of the formal languages used today were acquired. The most widespread is doubtless that hammered out by Russell and Whitehead in their great work ``Principia Mathematica'', which happily links the precision of Frege and the convenience of Peano {[}266{]}. Most actual formal languages are differentiated from it only by changes of secondary importance, aimed at simplifying its use. Among the most ingenious, we quote the ``functional'' writing of relations (for instance \(\in xy\) instead of \(x \in y\) , thought up by Lukasiewicz, thanks to which brackets can be completely omitted; but the most interesting is without doubt the introduction by Hilbert of the symbol \(\tau\) , which allows the consideration of the quantifiers \(\exists\) and \(\forall\) as abbreviation signs, the avoidance of the introduction of the ``universal'' functional symbol \(\iota\) of Peano and Russell (which is only applied to functional relations), and finally which avoids the need to formulate the axiom of choice in the theory of sets ({[}163 a{]}, v.III, p.~183).

\section*{THE NOTION OF TRUTH IN MATHEMATICS.}
\phantomsection
\addcontentsline{toc}{section}{THE NOTION OF TRUTH IN MATHEMATICS.}

Mathematicians have always been sure that they prove ``truths'' or ``true propositions''; such a conviction can obviously only be sentimental or metaphysical, and it is not by getting on to mathematical ground that it can be justified, nor even given a meaning that does not make it a tautology. The history of the concept of truth in mathematics is the concern therefore of the history of philosophy and not of that of mathematics; but the evolution of this concept has had an undeniable influence on that of mathematics, and because of this we can not let it go by in silence.

Let us note first that it is as rare to see a mathematician in possession of a strong philosophical culture as to see a philosopher who has an extensive knowledge of mathematics; the opinions of mathematicians on topics in philosophy, even when these questions are concerned with their field, are most often opinions received at second or third hand, coming from doubtful sources. But, precisely because of this, it is these average opinions which interest the mathematical historian, at least as much as the original views of thinkers such as Descartes or Leibniz (to mention two who were also mathematicians of the first rank), Plato (who at least kept up with the mathematics of his time), Aristotle or Kant (of whom the same could not be said).

The traditional notion of mathematical truth is that which goes back to the Renaissance. In this concept, there is no great difference between those objects which are the concern of mathematicians and those which are the concern of natural science; both are knowable and man has obtained a grasp on them both through intuition and through reason; there are no grounds to doubt either intuition or reason, which are fallible only if used incorrectly. ``One must have'', says Pascal, ``altogether a wrong spirit to reason badly about principles so great that it is almost impossible for them to escape'' ({[}244{]}, v. XII, p.~9). Descartes, by his fireside, convinced himself that ``there have been only Mathematicians who were able to find some proofs, that is to say some sure and certain reasons'' ({[}85 a{]}, v.VI, p.~19) and that (if his tale is to be believed) well before he had built up a metaphysics in which ``that very same thing'', he says, ``that I previously took as a rule, namely that the objects that we can visualise very clearly and very distinctly are all true, is only ensured because God is or exists, and that he is a perfect being'' ({[}85 a{]}, v. VI, p.~38). If Leibniz objects to Descartes that it cannot be seen how to recognise that an idea is ``clear and distinct'', \(^{23}\) he also considers axioms as obvious and undisputed consequences of the definitions as soon as the expressions are understood. \(^{24}\) It must not be forgotten either that, in the language of that time, mathematics includes many sciences that we do not recognise anymore as such, and sometimes even as far as the art of the engineer; and in the confidence that they inspire, the surprising success of their applications to ``natural philosophy'', to the ``mechanical arts'', to navigation, comes in for a large part.

In this way of looking at things, axioms are no more susceptible to being discussed or put into question than the rules of deduction; at most the choice can be left to each person, according to his preferences, to reason ``in the way of the ancients'' or to give free reign to his intuition. The choice of point of departure is also a matter of individual preference, and one sees appearing numerous ``editions'' of Euclid where the solid logical framework of the Elements becomes a strange travesty; some surveys are given of the infinitesimal calculus and rational mechanics, supposedly deduced from fundamentals which are remarkably badly established; and Spinoza was perhaps giving in good faith his Ethics as being proved in the manner of geometers ``more geometrico demonstrata''. If it is hard to find in the XVIIth century two mathematicians who agree on any matter whatsoever, if the polemics occur daily, endless and acrimonious, the notion of truth remains none the less not in question. ``Only having one truth about each object'', says Descartes, ``whoever finds it knows as much as can be known about it'' ({[}85 a{]}, v. VI, p.~21).

Although no Greek mathematical text of the high period on these questions has survived, it is probable that the point of view of Greek mathematicians on this subject had many more nuances. It is by experience only that the rules of deduction were elaborated to the point of inspiring complete confidence; before they could be considered to be above all discussion, it was necessary to go through mainly fumblings and paralogisms. It would also be to misconstrue the critical spirit of the Greeks, their taste for discussion and sophistry, to imagine that the ``axioms'' themselves that Pascal judged to be the most obvious (and that, according to a legend spread by his sister, he had, with infallible instinct, discovered himself during his childhood) were not the subject of long discussions. In an area that was not that of geometry strictly speaking, the paradoxes of the Eleates have preserved for us some traces of these polemics; and Archimedes when he makes the observation ({[}5 b{]}, v. II, p.~265) that his predecessors made use in many situations of the axiom to which we usually ascribe his name, adds that what is proved by means of this axiom ``has been accepted no less than what has been proved without it'', and that it is sufficient that his own results be accepted on the same basis. Plato, in accordance with his metaphysical views, presents mathematics as a means of access to a ``truth in itself'' and the objects it deals with as having a real existence in the world of ideas; he characterises none the less the mathematical method precisely in a famous extract from the Republic: ``Those that are involved with geometry and arithmetic \ldots{} assume even and odd, three kinds of angles; they treat them as known objects: once assumed, they esteem that they no longer need to give account of them either to themselves or to others, {[}considering it{]} as clear to each one; and starting there, they proceed in sequence, in order to reach by common agreement the goal that their research had suggested'' ({[}250{]}, Book VI, 510 c-e). That which makes up a proof is therefore firstly a point of departure supplying an arbitrary start (even though ``clear to everyone''), and beyond which, he says a bit further on, one does not try to dig; and then, a discussion that follows in sequence a series of intermediate stages; finally, at each step, the consent of the interlocutor guaranteeing the correctness of the reasoning. It must be added that once the axioms are stated, no new appeal to intuition is allowed on principle: Proclus, quoting Géminus, recalls that ``we have learnt from the pioneers of this science themselves, to take no account of conclusions which are only plausible when it is a case of reasoning which must be part of our geometric doctrine'' ({[}153 e{]}, v.I, p.~203).

Thus it is to experience and the cauldron of criticism that is due the elaboration of the rules of mathematical reasoning; and if it is true, as has been argued in a plausible way \([317\ d]\) , that Book VIII of Euclid has kept for us part of the arithmetic of Archytas, it is not surprising to see there the rigidity of the rather pedantic reasoning that does not fail to appear in all mathematical schools where ``rigour'' is discovered or believed to have been discovered. But, once having entered into the practice of mathematicians, it does not seem that these rules of reasoning have ever been doubted until quite recently: if with Aristotle and the Stoics, some of the rules are deduced from others by schemes of reasoning, the primitive rules are always assumed to be evident. In the same way, having gone back to the ``hypotheses'', ``axioms'' and ``postulates'' that appeared to them to supply a solid foundation for the science of their time (such as those for example that they put forward in the first ``Elements'' that tradition ascribes to Hippocrates of Chio, about 450 B. C.), the Greek mathematicians from the classical period seem to have bent all their efforts to the discovery of new results rather than to a critique of these foundations that, at that time, could not have failed to be sterile; and, putting aside all metaphysical preoccupations, it is from this general accord amongst mathematicians on the bases of their science that the text of Plato quoted above bears testimony.

On the other hand Greek mathematicians do not seem to have believed it possible to explain the ``primary notions'' that were their points of departure, straight line, surface, ratio of quantities; if they give ``definitions'' of them, it is obviously a crisis of conscience and without any illusions about their range. It goes without saying that on the other hand, about definitions other than those of ``primary notions'' (definitions often called ``nominal''), the Greek mathematicians and philosophers had perfectly clear ideas. It is in this context that intervenes explicitly, for the first time no doubt, the question of ``existence'' in mathematics. Aristotle does not miss observing that a definition does not imply the existence of the object defined, and that there must be over and above that either a postulate or a proof. No doubt his observation was derived from the practice of mathematicians; in any case Euclid takes care to postulate the existence of a circle, and to prove that of an equilateral triangle, of parallels, of a square, etc. as he introduces them into his arguments ({[}153 e{]}, Book I); these proofs are ``constructions''; in other words, he exhibits, basing himself on the axioms, mathematical objects that he proves satisfy the definitions that he needs to justify.

Thus we see Greek mathematics in the classical era reaching a kind of empirical certainty (whatever the metaphysical bases for it from some philosopher or other might be); if it cannot be conceived that the rules of reasoning can be questioned, the success of Greek science, and the feeling that exists of the inopportunity of a critical revision, play a great part in the confidence that the axioms themselves inspire, a confidence that is rather of the order of that (almost unlimited, as well) that was put during the last century in the principles of theoretical physics. It is in any case what is suggested by the motto of the school ``nihil est in intellectu quod non prius fuerit in sensu'', which is precisely that against which Descartes stands out, as not giving a firm enough basis to that which Descartes hoped to extract from the use of reason.

One must come down to the beginning of the XIXth century to see mathematicians recover from the arrogance of a Descartes (not to say that of a Kant or that of a Hegel, this latter somewhat late, as was fashionable, about the science of his time \(^{25}\) ), to a position as well balanced as that of the Greeks. The first blow to the classical concepts is the establishment of non-Euclidean hyperbolic geometry by Gauss, Lobatschevsky and Bolyai at the beginning of the century. We will not undertake to retrace here in detail the origins of this discovery, the outcome of numerous unfruitful attempts to prove the parallel postulate (see {[}105 a and b{]}). At the time its effect on the principles of mathematics was perhaps not so deep as is sometimes said. It simply forces the abandonment of the pretensions of the previous century to the ``absolute truth'' of Euclidean geometry, and all the more, the Leibnizian point of view that definitions imply axioms; these latter no longer appear as at all ``obvious'', but rather as hypotheses of which it must be determined whether they are adapted to the mathematical representation of the world of the senses. Gauss and Lobatschevsky believe that the debate between the different possible geometries can be determined by experience ({[}206{]}, p.~76). It is also the point of view of Riemann, of whom the famous inaugural Lecture ``On the hypotheses which form the foundation of geometry'' had as aim to furnish a general mathematical framework for the various natural phenomena: ``What remains to be resolved'' he says, ``is the question of knowing to what extent and up to what point these hypotheses are found to be confirmed by experience'' ({[}259 a{]}, p.~284). But that is a problem which visibly has nothing anymore to do with Mathematics; and none of the previous authors seem to put into question that, even if a ``geometry'' does not correspond to an experimental reality, its theorems remain no less ``mathematical truths''. \(^{26}\)

All the same, if it so, it is certainly no longer to an unlimited confidence in classical ``geometric intuition'' that such a conviction is due; the description that Riemann seeks to give of ``multiplicities n times extended'', the object of his work, only relies on ``intuitive'' \(^{27}\) considerations to reach a justification for the introduction of ``local co-ordinates''; from that moment on he apparently seems to feel himself on solid ground, namely that of Analysis. But this latter is grounded in the end on the concept of real number, which remained until then of a very intuitive nature; and progress in the theory of functions was leading in this respect to some very worrying results: with the research of Riemann himself on integration, and more so with the examples of curves without tangents, constructed by Bolzano and Weierstrass, it is the whole of the pathology of mathematics that was beginning. For a century we have seen so many monsters of this species that we are a bit blasé, and the most weird teratological characters must be accumulated in order still to astound us. But the effect produced on the majority of mathematicians of the XIXth century went from disgust to consternation: ``How'' H. Poincaré asks himself, ``can intuition deceive us at this point?'' ({[}251 d{]}, p.~19); and Hermite (not without a spark of humour which the commentators of this famous sentence do not all seem to have perceived) declares that he ``turns away with fear and horror from this lamentable plague of continuous functions that do not have a derivative'' ({[}160{]}, v. II, p.~318). The worst was that these phenomena, so contrary to common sense, could no longer be laid at the door of ideas badly elucidated, as at the time of the ``indivisibles'' (see p.~173), since they survived after the reform of Bolzano, Abel and Cauchy, who had established the foundation of the notion of limit in a manner as rigorous as the theory of ratios (see p.~154). It is thus fully to the gross and incomplete character of our geometric intuition that the account must be laid, and it is reasonable that since then it has remained discredited quite justifiably as a means of proof.

This realisation would inevitably react on classical mathematics, starting with geometry. In whatever respect the axiomatic construction of Euclid had been held, more than one imperfection had been noticed, and that already in antiquity. It was the postulate on parallels that had been the object of the greatest number of criticisms and attempts at proof; but the followers and commentators of Euclid had also attempted to prove other postulates (notably that of the equality of right angles) or recognised the insufficiency of certain definitions, such as those of the straight line or plane. In the XVIth century, Clavius, an editor of the Elements, notes the absence of a postulate guaranteeing the existence of the fourth proportional; for his part, Leibniz remarks that Euclid uses geometric intuition without mentioning it explicitly, for example when he admits (Elements, Book I, prop. 1) that two circles, each of which goes through the centre of the other have a common point ({[}198 b{]}, v. VII, p.~166). Gauss (who himself did not deny himself the use of such topological considerations) draws attention to the role played in Euclidean constructions by the notion of a point (or a straight line) being ``between'' two others, notion that is however not defined ({[}124 a{]}, v. VIII, p.~222). Finally, the use of displacements --- notably in the ``case of equality of triangles'' --- long assumed to be obvious, \(^{28}\) was soon to appear to the critics of the XIXth century as relying also on unstated axioms. One ends up thus, in the period from 1860 to 1885, with different partial revisions of the beginnings of geometry (Helmholtz, Méray, Houël) tending to remedy some of these gaps. But it is only with M. Pasch {[}245{]} that the abandonment of all appeals to intuition is a programme properly formulated and followed with full rigour. The success of his enterprise soon brought him numerous emulators who, principally between 1890 and 1910, gave quite varied statements of the axioms of Euclidean geometry. The most famous of these works were those of Peano, written in his symbolic language {[}246 d{]}, and especially the ``Grundlagen der Geometrie'' of Hilbert {[}163 c{]}, appearing in 1899, a book which, by its lucidity and the depth of its exposition, was to become immediately, with full justification, the charter for modern axiomatics, even to the extent of leading to forgetting its forerunners. It is indeed that, not content with giving there a complete system of axioms for Euclidean geometry, Hilbert classifies these axioms into different groups of different types, and sets himself to determine the exact area of influence of each of these groups of axioms, not only in developing the logical consequences of each of them in isolation, but also in discussing the different ``geometries'' obtained when one omits or modifies certain of these axioms (geometries amongst which those of Lobatschevsky and of Riemann appear only as particular cases \(^{29}\) ); he thus puts clearly in the picture, in an area considered until then as one of those nearest the reality of the senses, the freedom of which the mathematician disposes in his choice of postulates. In spite of the disarray caused in more than one philosopher by these ``metageometries'' with strange properties, the thesis of the ``Grundlagen'' was rapidly adopted almost unanimously by mathematicians; H. Poincaré, although hardly guilty of favouring formalism, recognised in 1902 that the axioms of geometry are conventions, for which the notion of ``truth'', as it is normally understood, has no more meaning ({[}251 c{]}, pp.~66-67). ``Mathematical truth'' resides thus uniquely in logical deduction starting from premises arbitrarily set by axioms. As will be seen later (pp.~35 ff.), the validity of the rules of reasoning used in making these deductions was itself soon to be put in question, bringing about thus a complete reshaping of the basic conceptions of mathematics.

\section*{OBJECTS, MODELS AND STRUCTURES.}
\phantomsection
\addcontentsline{toc}{section}{OBJECTS, MODELS AND STRUCTURES.}

\begin{enumerate}
\def\labelenumi{\Alph{enumi})}
\tightlist
\item
  Objects and structures of mathematics. --- From Antiquity to the XIXth century, there is a common understanding as to what are the principal objects of the mathematicians; they are exactly those that are mentioned by Plato in the passage quoted earlier (p.~13); numbers, quantities and figures. If, at first, must be added the objects and phenomena which are the subject of Mechanics, Astronomy, Optics and Music, these ``mathematical'' disciplines are always clearly separated, amongst the Greeks, from Arithmetic and Geometry, and starting with the Renaissance they fairly quickly attain the rank of independent sciences.
\end{enumerate}

Whatever the philosophical nuances by which the conception of mathematical objects are coloured by such and such a mathematician or philosopher, there is at least one point on which there is unanimity: it is that these objects are given to us and it is not in our power to assign to them arbitrary properties, in the same way as a physicist cannot change a natural phenomenon. Truth to say, psychological reactions doubtless form one part of these views, reactions which are not for us to follow up, but which each mathematician knows when he wears himself out in vain efforts to seize upon a proof which seems to slip away endlessly. From there to assimilate this resistance to the obstacles that the world of senses sets in our path there is only a single step; and even today, more than one, who claims an intransigent formalism, would subscribe voluntarily, in his innermost being, to this admission of Hermite: ``I believe that numbers and the functions of Analysis are not the arbitrary result of our minds; I think that they exist outside of us, with the same character of necessity as the things of objective reality, and we meet them or discover them, and study them, as do the physicists, the chemists and the zoologists'' ({[}160{]}, v. II, p.~398).

There is no question in the classical conception of mathematics, of straying away from the study of numbers and figures; but this official doctrine, to which every mathematician believes himself bound to bring his verbal adherence, becomes no less bit by bit an intolerable burden, as new ideas accumulate. The embarrassment of the algebraists up against negative numbers does not cease until analytical Geometry gives it a convenient ``interpretation''; but, well into the XVIIIth century still, d' Alembert, discussing the question in the Encyclopédie ({[}75 a{]}, article NÉGATIF), loses his nerve suddenly after a column of fairly confused explanations, and is content to conclude that ``the rules of algebraic operations on negative quantities are generally assumed by everybody and perceived generally as exact, whatever idea is linked elsewhere to these quantities''. As for imaginary numbers, the scandal is far bigger; for if they are ``impossible'' roots and if (until about 1800) no way of ``interpreting'' them is seen, how can one, without contradiction, talk of these indefinable beings, and above all why introduce them? D'Alembert keeps here a prudent silence and does not even state these questions, no doubt because he realises that he would not be able to answer other than was done naively by A. Girard a century earlier ({[}129{]}, f.~22): ``It could be said: what use are these impossible solutions? I answer: for three things, for the certainty of the general rule, and that there is no other solution, and for its usefulness.''

In Analysis, the situation, in the XVIIth century, was no better. It is a happy circumstance that analytical Geometry was to appear, as if at a designated point, to give a ``representation'' in the shape of a geometric figure, of the great creation of the XVIIth century, the notion of a function, and so assist powerfully (with Fermat, Pascal or Barrow) at the birth of the infinitesimal Calculus (cf.~p.~193). But it is known on the other hand, to what philosophico-mathematical controversies the notions of the infinitely small and indivisible were going to give rise. And if d'Alembert is happier here, and recognises that in the ``metaphysics'' of the infinitesimal Calculus there is nothing other than the notion of limit ({[}75 a{]}, articles DIFFÉRENTIEL and LIMITE, and {[}75 b{]}), he is no more able than his contemporaries, to understand the real meaning of expansion in divergent series, and to explain the paradox of exact results obtained at the end of calculations with expressions deprived of any numerical interpretation. Finally, even in the domain of ``geometric certainty'', the Euclidean framework bursts: when Stirling, in 1717, does not shrink from saying that a certain curve has a ``double imaginary point at infinity'' ({[}299{]}, p.~93 of the new edition), he would certainly be in trouble trying to link such an ``object'' to commonly understood notions; and Poncelet who, at the beginning of the XIXth century, gives considerable development to such ideas in founding projective geometry (see p.~131), is content still to invoke as justification a ``principle of continuity'' which is altogether metaphysical.

One can conceive that, under these conditions (and at the exact moment when, paradoxically, the ``absolute truth'' of mathematics is proclaimed with the greatest force), the notion of proof seems to become more and more blurred during the XVIIIth century, since it is out of the question to set down, as did the Greeks, the notions on which reasoning is conducted, and their fundamental properties. The return towards rigour, which starts at the beginning of the XIXth century, brings some improvement to this state of things, but does not stop for all that the stream of new notions: one sees thus appearing in Algebra the imaginaries of Galois ({[}123{]}, p.~113-127), the ideal numbers of Kummer {[}188 b{]}, which are followed by vectors and quaternions, n-dimensional spaces, multivectors and tensors (see pp.~61 ff.), not to speak of Boolean algebra. No doubt great progress (which precisely allows the return to rigour, without losing any of the conquests of previous times) occurs with the possibility of giving ``models'' for these new notions in more classical terms: the ideal numbers or the imaginaries of Galois are interpreted by means of the theory of congruences (see pp.~82 ff.), n-dimensional geometry only appears (if so desired) as a pure language for stating results in algebra ``with n variables''; and for the classical imaginary numbers --- of which the geometric representation by the points of a plane (see pp.~161 ff.) marks the beginning of this flowering of Algebra --- there is soon the choice between the geometric ``model'' and an interpretation in terms of congruences (cf.~p.~82). But mathematicians begin at last to feel sharply that it is to fight against a natural slope along which their work is dragging them, and that it must be allowable, in mathematics, to reason about objects that have no sensible ``interpretation'': ``It is not of the essence of mathematics'', says Boole in 1854, ``to be occupied with the ideas of number and quantity'' ({[}29{]}, v. II, p.~13). \(^{30}\) The same preoccupation leads Grassmann, in his ``Ausdehnungslehre''

of 1844, to present his calculus in a form from which the notions of number or of geometric object are, at first, excluded. \(^{31}\) And, a bit later, Riemann, in his inaugural Lecture, takes care at the beginning not to speak of ``points'', but rather of ``determinations'' (Bestimmungsweise), in his description of ``multiplicities n times extended'', and underlines that, in such a multiplicity, the ``metric relations'' (Massverhältnisse) ``can only be studied for abstract quantities and can only be represented by formulae; under certain conditions, one can however decompose them into relations of which each taken in isolation is susceptible to a geometric representation, and in that way it is possible to express the results of the calculation in a geometric form'' ({[}259 a{]}, p.~276).

From that moment, the expansion of the axiomatic method is an accomplished fact. If, for a while longer, it is believed useful to control, when possible, the ``abstract'' results by geometric intuition, at least it is admitted that the ``classical'' objects are no longer the only ones that the mathematician can legitimately study. It is that --- precisely because of the multiple ``interpretations'' or ``models'' that are possible --- it has been recognised that the ``nature'' of mathematical objects is in the end secondary, and that it matters little, for example, that a result is presented as a theorem of ``pure'' geometry, or as a theorem of algebra by the means of analytical geometry. In other words, the essence of mathematics --- this elusive notion that could until then only have been expressed by vague names such as ``general rule'' or ``metaphysics'' --- appears as the study of relationships between objects that are only (voluntarily) known and described by some of their properties, precisely those that are put as axioms at the foundations of their theory. It is this that had already been clearly seen by Boole in 1847, when he wrote that mathematics deals with ``operations considered in themselves, independently of the diverse objects to which they can be applied'' ({[}29{]}, v. I, p.~3). Hankel, in 1867, inaugurating the axiomatisation of algebra, defends a mathematics that is ``purely intellectual, a pure theory of forms, which has as its purpose, not the combination of quantities, or of their images, the numbers, but objects of thought (``Gedankendinge'') to which may correspond effective objects or relations, even though such a correspondence is not necessary'' ({[}146{]}, p.10). Cantor, in 1883, echoes this claim of a ``free mathematics'' by proclaiming that ``mathematics is entirely free in its development, and its concepts are only linked by the necessity of being consistent, and are co-ordinated with concepts introduced previously by means of precise definitions'' ({[}47{]}, p.182). Finally, the revision of Euclidean geometry succeeds in spreading and popularising these ideas. Pasch himself, although still attached to a certain ``reality'' of geometric objects, recognises that geometry is in fact independent of their significance, and consists purely in the study of their relations ({[}245{]}, p.90); a concept that Hilbert pushes to its logical conclusion in underlining that even the names of the basic notions of a mathematical theory can be chosen arbitrarily, \(^{32}\) and that Poincaré expresses by saying that the axioms are ``disguised definitions'', thus completely reversing the scholastic point of view.

It would thus be tempting to say that the modern notion of ``structure'' is attained in substance around 1900; in fact it will need still another thirty years of apprenticeship before it appears in all its glory. It is no doubt not difficult to recognise structures of the same kind when they have a fairly simple nature; for the group structure, for example, this stage was reached already by the middle of the XIXth century. But at the same time, Hankel can be seen fighting --- without quite succeeding --- to bring out the general ideas of field and extension, that he does not succeed in expressing except in the form of a ``principle of permanence'' which is semi metaphysical {[}146{]}, and which will be only formulated in a definitive form by Steinitz {[}294 a{]} 40 years later. Above all it has been fairly difficult, in this matter, to free oneself from the impression that mathematical objects are ``given'' to us with their structure; only a fairly long use of functional Analysis has been able to familiarise modern mathematicians with the idea that, for example, there are several ``natural'' topologies on the rational numbers, and several measures on the number line. With this disassociation the transition to general definition of structures was finally realised.

\begin{enumerate}
\def\labelenumi{\Alph{enumi})}
\setcounter{enumi}{1}
\item
  Models and isomorphisms. --- The intervention of the notion of ``model'' or ``interpretation'' of one mathematical theory by means of another will have been noticed on several occasions. That is not a recent idea, and no doubt what can be seen here is a manifestation ceaselessly recurring of a deep feeling for the unity of the various ``mathematical sciences''. If one can take as authentic the traditional maxim ``All things are numbers'' of the first Pythagoreans, that can be considered as the trace of a first attempt to bring back the geometry and algebra of the times to arithmetic. Although the discovery of the irrationals seemed to close for ever that route, the reaction to which it gave birth in Greek mathematics was a second attempt at synthesis this time taking geometry as the basis, and absorbing among others the solution methods for algebraic equations inherited from the Babylonians. \(^{33}\) It is known that this conception was to remain until the fundamental reform of R. Bombelli and of Descartes, assimilating every measure of quantity to a measure of length (in other words, to a real number; cf.~p.~151). But with the creation of analytical geometry by Descartes and Fermat, the tendency is again reversed, and a much tighter fusion of geometry and algebra is obtained, but this time to the benefit of algebra. Further, elsewhere, Descartes goes further and conceives of the essential unity of ``all these sciences that are commonly called Mathematics \ldots{} Even though their subjects are different'', he says ``they still do not fail to converge, in that they consider nothing other than the various relations or ratios that are to be found there'' ({[}85 a{]}, v. VI, pp.~19-20). \(^{34}\) However, this point of view only tended to make Algebra the fundamental mathematical science; a conclusion against which Leibniz protests vigorously, who also himself, as has been seen, conceived of a ``universal Mathematics'', but on a much vaster scale and already quite close to modern ideas. Making precise the ``accord'' of which Descartes spoke, he glimpses, in fact, for the first time, the general notion of isomorphism (which he calls ``similitude''), and the possibility of ``identifying'' relations or operations that are isomorphic; he gives examples addition and multiplication ({[}69 a{]}, pp.~301-303). But these audacious views remained without echo amongst his contemporaries, and one must await the expansion of Algebra which takes place around the middle of the XIXth century (see pp.~51 ff.) to see the beginnings of the realisation of the Leibnizian dreams. We have already underlined that it is at that moment that the ``models'' multiply and that it becomes usual to go from one theory to another by a simple change of language; the most striking example of this is perhaps duality in projective geometry (see p.~132), where the practice, frequent at that time, of printing face to face, in two columns, the theorems ``dual'' one to the other, doubtless plays a great part in the full realisation of the notion of isomorphism. From a more technical point of view, certainly the notion of isomorphic groups for abelian groups is known to Gauss, for groups of permutations to Galois (cf.~pp.~51 ff.); it is acquired in a general way for any group whatsoever around the middle of the XIXth century. \(^{35}\) Following this, with each new axiomatic theory, it is a natural development to define a notion of isomorphism; but it is only with the modern notion of structure that it was finally recognised that every structure carries within itself a notion of isomorphism, and that it is not necessary to give a special definition of it for each type of structure.
\item
  The arithmetisation of classical mathematics. --- The use, more and more widespread, of the notion of ``model'' was also going to allow the realisation in the XIXth century of the unification of mathematics dreamed of by the Pythagoreans. At the beginning of the century, whole numbers and continuous quantities seemed as irreconcilable as ever in antiquity; real numbers remain linked to the notion of geometric quantity (at least to that of length), and it is to this latter that appeal had been made for the ``models'' of negative numbers and imaginary numbers. Even rational numbers were traditionally attached to the idea of the splitting of a quantity into equal parts; only the whole numbers remained apart, as ``exclusively the product of our minds'' as is said by Gauss in 1832, putting them in opposition to the notion of space ({[}124 a{]}, v. VIII, p.~201). The first efforts to bring together Arithmetic and Analysis were made at first with rational numbers (positive and negative) and are due to Martin Ohm (1822); they were taken up again around 1860 by several authors, notably Grassmann, Hankel and Weierstrass (in his unpublished lectures); it is to this latter that it appears is due the idea of obtaining a ``model'' of the positive rational numbers or of whole negative numbers by considering classes of pairs of whole numbers. But the most important step remained to be taken, namely to find a ``model'' of irrational numbers in the theory of rational numbers; around 1870, it had become an urgent problem, in view of the necessity, after the discovery of ``pathological'' phenomena in Analysis, to eliminate every trace of geometric intuition in the vague notion of ``quantity'' in the definition of real numbers. It is known that this problem was solved around this time, almost simultaneously by Cantor, Dedekind, Méray and Weierstrass, and using fairly different methods (see p.~155).
\end{enumerate}

From this moment, whole numbers became the foundation of all classical mathematics. Further, ``models'' based on Arithmetic acquired still more importance with the extension of the axiomatic method and the conception of mathematical objects as free creations of the mind. There remained in fact one restriction to this freedom claimed by Cantor, the question of ``existence'' which had already preoccupied the Greeks, and which arose here in a much more immediate way, precisely since all appeals to an intuitive representation had now been abandoned. We will see later (pp.~36 ff.) of what a philosophico-mathematical maelstrom the notion of ``existence'' was going to be the centre in the first years of the XXth century. But in the XIXth century that had not been reached, and to prove the existence of a mathematical object having given properties, it was simply, as for Euclid, ``constructing'' an object with given properties. That was precisely what arithmetic ``models'' were for: once the real numbers were ``interpreted'' in terms of whole numbers, complex numbers and Euclidean geometry were also, thanks to analytical Geometry, and it was the same for all new algebraic objects introduced since the beginning of the century; finally --- a discovery that had a great effect --- Beltrami and Klein had even obtained Euclidean ``models'' of the non-Euclidean geometries of Lobatschevsky and of Riemann (see p.~134), and in consequence ``arithmetised'' (and thereby completely justified) these theories which had at first sight aroused such distrust.

\begin{enumerate}
\def\labelenumi{\Alph{enumi})}
\setcounter{enumi}{3}
\tightlist
\item
  The axiomatisation of arithmetic.---It was in line with this evolution that a subsequent turn towards the foundations of arithmetic itself was made, and indeed that is what can be seen around 1880. It appears that before the XIXth century no attempt had been made to define addition and multiplication of whole numbers except by a direct appeal to intuition; Leibniz is the only one who, faithful to his principles, warns expressly that ``truths'' as ``obvious'' as \(2 + 2 = 4\) are no less susceptible to proof if one thinks about the definitions of the numbers involved ({[}198 b{]}, v. IV, p.~403; cf.~{[}69 a{]}, p.~203); and he did not consider at all that commutativity of addition and multiplication were intrinsic. \(^{36}\) But he does not take his thoughts on this subject any further, and about the middle of the XIXth century, still no progress in this direction had been made: Weierstrass himself, whose lectures contributed considerably to extend the ``arithmeticising'' point of view, does not seem to have felt the need for a logical clarification of the theory of whole numbers. The first steps in this direction seem to be due to Grassmann, who, in 1861 ({[}134{]}, v. II \(_{2}\) , p.~295) gives a definition of addition and multiplication for whole numbers, and proves their fundamental properties (commutativity, associativity, distributativity) using only the operation \(x \mapsto x + 1\) and the principle of induction. This latter had been clearly conceived and used for the first time in the XVIIth century by B. Pascal ({[}244{]}, v. III, p.~456) \(^{37}\) --- even though one can find in Antiquity some more or less conscious applications of it --- and was currently used by mathematicians since the second half of the XVIIth century. But it is only in 1888 that Dedekind ({[}79{]}, v. III, pp.359-361) stated a complete system of axioms for arithmetic (a system repeated 3 years later by Peano and usually known by his name {[}246 c{]}), which contained in particular a precise formulation of the principle of induction (that Grassmann uses yet without stating it explicitly).
\end{enumerate}

With this axiomatisation, it seemed that the definitive foundations of mathematics had been reached. In fact, at the very moment when the axioms of arithmetic were being clearly formulated, these same, for many mathematicians (starting with Dedekind and Peano themselves) were already deprived of this role of primordial science, in favour of the latest arrival among the theories of mathematics, the theory of sets; and the controversies that were going to unravel around the notion of whole number cannot be isolated from the great ``crisis of the foundations'' of the years 1900-1930.
